\section{BPMN Support And Limits}
\label{sec:SupportLimits}

Durga supports a pragmatic BPMN subset that is broad enough for exploration and code
generation, but not yet a complete BPMN execution engine.

\subsection{Well-supported areas}

The strongest current coverage includes:
\begin{itemize}
\item start events and explicit end completion handling,
\item service, user and manual task mappings,
\item XOR and parallel gateway routing,
\item intermediate timer, message and signal events,
\item task and subprocess boundary timer, error and escalation handling,
\item call activities,
\item embedded subprocesses including nested subprocesses,
\item event subprocesses for message, signal, timer, error and escalation starts.
\end{itemize}

\subsection{Current semantic boundaries}

Some behaviors are supported only to the level needed for explicit generated runtimes, not
to full BPMN-engine fidelity:
\begin{itemize}
\item cancellation propagation is pragmatic rather than globally token-precise,
\item subprocess scope handling is explicit but still lighter than a full BPMN execution scope manager,
\item monitoring is strong on current state, lighter on long-term history and trend analysis,
\item generated runtimes are optimized for understanding and experimentation rather than minimalism.
\end{itemize}

\subsection{Still outside the current focus}

The main areas still outside the project's current center of gravity are:
\begin{itemize}
\item full production-grade workflow-engine semantics,
\item richer compensation and advanced BPMN exception orchestration,
\item large-scale operational hardening beyond local and exploratory usage,
\item a generalized multi-node monitoring and query deployment model.
\end{itemize}

The repository coverage matrix remains the most detailed element-by-element reference, but
the overall design intent is simple: support enough BPMN to make the Kafka mapping useful,
while keeping the generated runtime explicit and teachable.
